
% https://akaedu.github.io/book/ch31.html

\chapter{Shell脚本}
\BookAnchor{db.ch31}\BookAnchor{db.ch31.id2871639}%


\BookNav{\hyperlink{db.ch30s05}{上一页}\enspace\textperiodcentered\enspace \hyperlink{db.pt03}{上一级}\enspace\textperiodcentered\enspace \hyperlink{db.index}{起始页}\enspace\textperiodcentered\enspace \hyperlink{db.ch31s01}{下一页}}


% https://akaedu.github.io/book/ch31s01.html

\section{Shell的历史}
\BookAnchor{db.ch31s01}\BookAnchor{db.ch31s01.id2871651}%
Shell的作用是解释执行用户的命令，用户输入一条命令，Shell就解释执行一条，这种方式称为交互式（Interactive）\BookAnchor{db.ch31s01.id2871030}，Shell还有一种执行命令的方式称为批处理（Batch）\BookAnchor{db.ch31s01.id2871037}，用户事先写一个Shell脚本（Script）\BookAnchor{db.ch31s01.id2871044}，其中有很多条命令，让Shell一次把这些命令执行完，而不必一条一条地敲命令。Shell脚本和编程语言很相似，也有变量和流程控制语句，但Shell脚本是解释执行的，不需要编译，Shell程序从脚本中一行一行读取并执行这些命令，相当于一个用户把脚本中的命令一行一行敲到Shell提示符下执行。



由于历史原因，UNIX系统上有很多种Shell：


\begin{enumerate}

\item \BookCode{sh}（Bourne Shell）：由Steve Bourne开发，各种UNIX系统都配有\BookCode{sh}。

\item \BookCode{csh}（C Shell）：由Bill Joy开发，随BSD UNIX发布，它的流程控制语句很像C语言，支持很多Bourne Shell所不支持的功能：作业控制，命令历史，命令行编辑。

\item \BookCode{ksh}（Korn Shell）：由David Korn开发，向后兼容\BookCode{sh}的功能，并且添加了\BookCode{csh}引入的新功能，是目前很多UNIX系统标准配置的Shell，在这些系统上\BookCode{/\BookSplit{}bin/\BookSplit{}sh}往往是指向\BookCode{/\BookSplit{}bin/\BookSplit{}ksh}的符号链接。

\item \BookCode{tcsh}（TENEX C Shell）：是\BookCode{csh}的增强版本，引入了命令补全等功能，在FreeBSD、Mac OS X等系统上替代了\BookCode{csh}。

\item \BookCode{bash}（Bourne Again Shell）：由GNU开发的Shell，主要目标是与POSIX标准保持一致，同时兼顾对\BookCode{sh}的兼容，\BookCode{bash}从\BookCode{csh}和\BookCode{ksh}借鉴了很多功能，是各种Linux发行版标准配置的Shell，在Linux系统上\BookCode{/\BookSplit{}bin/\BookSplit{}sh}往往是指向\BookCode{/\BookSplit{}bin/\BookSplit{}bash}的符号链接\BookAnchor{db.ch31s01.id2871814}\footnote[38]{\BookAnchor{db.ch31s01.ftn.id2871814}\BookFootnoteBack{db.ch31s01.id2871814}最新的发行版有一些变化，例如Ubuntu 7.10的\BookCode{/\BookSplit{}bin/\BookSplit{}sh}是指向\BookCode{/\BookSplit{}bin/\BookSplit{}dash}的符号链接，\BookCode{dash}也是一种类似\BookCode{bash}的Shell。
\BookTerminal{code/ch31s01/01.txt}{5}{3}}。虽然如此，\BookCode{bash}和\BookCode{sh}还是有很多不同的，一方面，\BookCode{bash}扩展了一些命令和参数，另一方面，\BookCode{bash}并不完全和\BookCode{sh}兼容，有些行为并不一致，所以\BookCode{bash}需要模拟\BookCode{sh}的行为：当我们通过\BookCode{sh}这个程序名启动\BookCode{bash}时，\BookCode{bash}可以假装自己是\BookCode{sh}，不认扩展的命令，并且行为与\BookCode{sh}保持一致。
\end{enumerate}


文件\BookCode{/\BookSplit{}etc/\BookSplit{}shells}给出了系统中所有已知（不一定已安装）的Shell，除了上面提到的Shell之外还有很多变种。


\BookPlain{code/ch31s01/02.txt}{9}{14}


用户的默认Shell设置在\BookCode{/\BookSplit{}etc/\BookSplit{}passwd}文件中，例如下面这行对用户mia的设置:


\BookPlain{code/ch31s01/03.txt}{3}{1}


用户mia从字符终端登录或者打开图形终端窗口时就会自动执行\BookCode{/\BookSplit{}bin/\BookSplit{}bash}。如果要切换到其它Shell，可以在命令行输入程序名，例如：


\BookTerminal{code/ch31s01/04.txt}{7}{5}


本章只介绍\BookCode{bash}和\BookCode{sh}的用法和相关语法，不介绍其它Shell。所以下文提到Shell都是指\BookCode{bash}或\BookCode{sh}。

\BookNav{\hyperlink{db.ch31}{上一页}\enspace\textperiodcentered\enspace \hyperlink{db.ch31}{上一级}\enspace\textperiodcentered\enspace \hyperlink{db.index}{起始页}\enspace\textperiodcentered\enspace \hyperlink{db.ch31s02}{下一页}}


% https://akaedu.github.io/book/ch31s02.html

\section{Shell如何执行命令}
\BookAnchor{db.ch31s02}\BookAnchor{db.ch31s02.id2872012}%
\subsection{执行交互式命令}
\BookAnchor{db.ch31s02.id2872017}%


用户在命令行输入命令后，一般情况下Shell会\BookCode{fork}并\BookCode{exec}该命令，但是Shell的内建命令例外，执行内建命令相当于调用Shell进程中的一个函数，并不创建新的进程。以前学过的\BookCode{cd}、\BookCode{alias}、\BookCode{umask}、\BookCode{exit}等命令即是内建命令，凡是用\BookCode{which}命令查不到程序文件所在位置的命令都是内建命令，内建命令没有单独的man手册，要在man手册中查看内建命令，应该


\BookTerminal{code/ch31s02/01.txt}{3}{1}


本节会介绍很多内建命令，如\BookCode{export}、\BookCode{shift}、\BookCode{if}、\BookCode{eval}、\BookCode{[\BookSplit{}}、\BookCode{for}、\BookCode{while}等等。内建命令虽然不创建新的进程，但也会有Exit Status，通常也用0表示成功非零表示失败，虽然内建命令不创建新的进程，但执行结束后也会有一个状态码，也可以用特殊变量\BookCode{\$?}读出。


\subsubsection*{习题}
\BookAnchor{db.ch31s02.id2872131}%


1、在完成\hyperlink{db.ch30s05.id2868354}{第 5 节 “练习：实现简单的Shell”}时也许有的读者已经试过了，在自己实现的Shell中不能执行\BookCode{cd}命令，因为\BookCode{cd}是一个内建命令，没有程序文件，不能用\BookCode{exec}执行。现在请完善该程序，实现\BookCode{cd}命令的功能，用\BookCode{chdir(\BookSplit{}2)\BookSplit{}}函数可以改变进程的当前工作目录。



2、思考一下，为什么\BookCode{cd}命令要实现成内建命令？可不可以实现一个独立的\BookCode{cd}程序，例如\BookCode{/\BookSplit{}bin/\BookSplit{}cd}，就像\BookCode{/\BookSplit{}bin/\BookSplit{}ls}一样？


\subsection{执行脚本}
\BookAnchor{db.ch31s02.id2872211}%


首先编写一个简单的脚本，保存为\BookCode{script.\BookSplit{}sh}：


\BookExampleTitle{简单的Shell脚本}\BookAnchor{db.ch31s02.id2872225}

\BookSource{numbers=none}{code/ch31s02/02.pv}{6}{4}


Shell脚本中用\BookCode{\#}表示注释，相当于C语言的\BookCode{/\BookSplit{}/\BookSplit{}}注释。但如果\BookCode{\#}位于第一行开头，并且是\BookCode{\#!}（称为Shebang\BookAnchor{db.ch31s02.id2872262}）则例外，它表示该脚本使用后面指定的解释器\BookCode{/\BookSplit{}bin/\BookSplit{}sh}解释执行。如果把这个脚本文件加上可执行权限然后执行：


\BookTerminal{code/ch31s02/03.txt}{4}{2}


Shell会\BookCode{fork}一个子进程并调用\BookCode{exec}执行\BookCode{.\BookSplit{}/\BookSplit{}script.\BookSplit{}sh}这个程序，\BookCode{exec}系统调用应该把子进程的代码段替换成\BookCode{.\BookSplit{}/\BookSplit{}script.\BookSplit{}sh}程序的代码段，并从它的\BookCode{\_\BookSplit{}start}开始执行。然而\BookCode{script.\BookSplit{}sh}是个文本文件，根本没有代码段和\BookCode{\_\BookSplit{}start}函数，怎么办呢？其实\BookCode{exec}还有另外一种机制，如果要执行的是一个文本文件，并且第一行用Shebang指定了解释器，则用解释器程序的代码段替换当前进程，并且从解释器的\BookCode{\_\BookSplit{}start}开始执行，而这个文本文件被当作命令行参数传给解释器。因此，执行上述脚本相当于执行程序


\BookTerminal{code/ch31s02/04.txt}{3}{1}


以这种方式执行不需要\BookCode{script.\BookSplit{}sh}文件具有可执行权限。再举个例子，比如某个\BookCode{sed}脚本的文件名是\BookCode{script}，它的开头是


\BookPlain{code/ch31s02/05.txt}{3}{1}


执行\BookCode{.\BookSplit{}/\BookSplit{}script}相当于执行程序


\BookTerminal{code/ch31s02/06.txt}{3}{1}


以上介绍了两种执行Shell脚本的方法：


\BookTerminal{code/ch31s02/07.txt}{4}{2}


这两种方法本质上是一样的，执行上述脚本的步骤为：


\BookFigure{Shell脚本的执行过程}{\BookImage{images/shellscript.shellexec.png}{0.902}}{\BookAnchor{db.ch31s02.id2872410}}

\begin{enumerate}

\item 交互Shell（\BookCode{bash}）\BookCode{fork}/\BookCode{exec}一个子Shell（\BookCode{sh}）用于执行脚本，父进程\BookCode{bash}等待子进程\BookCode{sh}终止。

\item \BookCode{sh}读取脚本中的\BookCode{cd .\BookSplit{}.\BookSplit{}}命令，调用相应的函数执行内建命令，改变当前工作目录为上一级目录。

\item \BookCode{sh}读取脚本中的\BookCode{ls}命令，\BookCode{fork}/\BookCode{exec}这个程序，列出当前工作目录下的文件，\BookCode{sh}等待\BookCode{ls}终止。

\item \BookCode{ls}终止后，\BookCode{sh}继续执行，读到脚本文件末尾，\BookCode{sh}终止。

\item \BookCode{sh}终止后，\BookCode{bash}继续执行，打印提示符等待用户输入。
\end{enumerate}


如果将命令行下输入的命令用()括号括起来，那么也会\BookCode{fork}出一个子Shell执行小括号中的命令，一行中可以输入由分号;隔开的多个命令，比如：


\BookTerminal{code/ch31s02/08.txt}{3}{1}


和上面两种方法执行Shell脚本的效果是相同的，\BookCode{cd .\BookSplit{}.\BookSplit{}}命令改变的是子Shell的\BookCode{PWD}，而不会影响到交互式Shell。然而命令


\BookTerminal{code/ch31s02/09.txt}{3}{1}


则有不同的效果，\BookCode{cd .\BookSplit{}.\BookSplit{}}命令是直接在交互式Shell下执行的，改变交互式Shell的\BookCode{PWD}，然而这种方式相当于这样执行Shell脚本：


\BookTerminal{code/ch31s02/10.txt}{3}{1}


或者


\BookTerminal{code/ch31s02/11.txt}{3}{1}


\BookCode{source}或者\BookCode{.\BookSplit{}}命令是Shell的内建命令，这种方式也不会创建子Shell，而是直接在交互式Shell下逐行执行脚本中的命令。


\subsubsection*{习题}
\BookAnchor{db.ch31s02.id2872645}%


1、解释如下命令的执行过程：


\BookTerminal{code/ch31s02/12.txt}{5}{3}

\BookNav{\hyperlink{db.ch31s01}{上一页}\enspace\textperiodcentered\enspace \hyperlink{db.ch31}{上一级}\enspace\textperiodcentered\enspace \hyperlink{db.index}{起始页}\enspace\textperiodcentered\enspace \hyperlink{db.ch31s03}{下一页}}


% https://akaedu.github.io/book/ch31s03.html

\section{Shell的基本语法}
\BookAnchor{db.ch31s03}\BookAnchor{db.ch31s03.id2872661}%
\subsection{变量}
\BookAnchor{db.ch31s03.id2872666}%


按照惯例，Shell变量由全大写字母加下划线组成，有两种类型的Shell变量：


\begin{description}
\item[{环境变量}]


在\hyperlink{db.ch30s02.id2865687}{第 2 节 “环境变量”}中讲过，环境变量可以从父进程传给子进程，因此Shell进程的环境变量可以从当前Shell进程传给\BookCode{fork}出来的子进程。用\BookCode{printenv}命令可以显示当前Shell进程的环境变量。

\item[{本地变量}]


只存在于当前Shell进程，用\BookCode{set}命令可以显示当前Shell进程中定义的所有变量（包括本地变量和环境变量）和函数。


\end{description}


环境变量是任何进程都有的概念，而本地变量是Shell特有的概念。在Shell中，环境变量和本地变量的定义和用法相似。在Shell中定义或赋值一个变量：


\BookTerminal{code/ch31s03/01.txt}{3}{1}


注意等号两边都不能有空格，否则会被Shell解释成命令和命令行参数。



一个变量定义后仅存在于当前Shell进程，它是本地变量，用\BookCode{export}命令可以把本地变量导出为环境变量，定义和导出环境变量通常可以一步完成：


\BookTerminal{code/ch31s03/02.txt}{3}{1}


也可以分两步完成：


\BookTerminal{code/ch31s03/03.txt}{4}{2}


用\BookCode{unset}命令可以删除已定义的环境变量或本地变量。


\BookTerminal{code/ch31s03/04.txt}{3}{1}


如果一个变量叫做\BookCode{VARNAME}，用\BookCode{\$\{VARNAME\}}可以表示它的值，在不引起歧义的情况下也可以用\BookCode{\$VARNAME}表示它的值。通过以下例子比较这两种表示法的不同：


\BookTerminal{code/ch31s03/05.txt}{6}{4}


注意，在定义变量时不用\$，取变量值时要用\$。和C语言不同的是，Shell变量不需要明确定义类型，事实上Shell变量的值都是字符串，比如我们定义\BookCode{VAR=\BookSplit{}45}，其实\BookCode{VAR}的值是字符串\BookCode{45}而非整数。Shell变量不需要先定义后使用，如果对一个没有定义的变量取值，则值为空字符串。


\subsection{文件名代换（Globbing）：* ? []}
\BookAnchor{db.ch31s03.id2872839}\BookAnchor{db.ch31s03.id2872843}%


这些用于匹配的字符称为通配符（Wildcard）\BookAnchor{db.ch31s03.id2872854}，具体如下：


\begingroup
\setlength{\tabcolsep}{6pt}
\small
\sloppy
\renewcommand{\arraystretch}{1.16}
\setlength{\parskip}{0pt}
\begin{longtable}{L{\dimexpr 0.264609\linewidth-2\tabcolsep\relax}L{\dimexpr 0.735391\linewidth-2\tabcolsep\relax}}
\BookTableTitle{通配符}\BookTableLabel{\BookAnchor{db.ch31s03.id2872862}}
\toprule
* & 匹配0个或多个任意字符 \\
\addlinespace[2pt]
? & 匹配一个任意字符 \\
\addlinespace[2pt]
[若干字符] & 匹配方括号中任意一个字符的一次出现 \\
\bottomrule
\end{longtable}
\endgroup

\BookTerminal{code/ch31s03/06.txt}{6}{4}


注意，Globbing所匹配的文件名是由Shell展开的，也就是说在参数还没传给程序之前已经展开了，比如上述\BookCode{ls ch0[\BookSplit{}012]\BookSplit{}.\BookSplit{}doc}命令，如果当前目录下有\BookCode{ch00.\BookSplit{}doc}和\BookCode{ch02.\BookSplit{}doc}，则传给\BookCode{ls}命令的参数实际上是这两个文件名，而不是一个匹配字符串。


\subsection{命令代换：`或 \$()}
\BookAnchor{db.ch31s03.id2872936}%


由反引号括起来的也是一条命令，Shell先执行该命令，然后将输出结果立刻代换到当前命令行中。例如定义一个变量存放\BookCode{date}命令的输出：


\BookTerminal{code/ch31s03/07.txt}{4}{2}


命令代换也可以用\BookCode{\$(\BookSplit{})\BookSplit{}}表示：


\BookTerminal{code/ch31s03/08.txt}{3}{1}

\subsection{算术代换：\$(())}
\BookAnchor{db.ch31s03.id2872971}%


用于算术计算，\BookCode{\$(\BookSplit{}(\BookSplit{})\BookSplit{})\BookSplit{}}中的Shell变量取值将转换成整数，例如：


\BookTerminal{code/ch31s03/09.txt}{4}{2}


\BookCode{\$(\BookSplit{}(\BookSplit{})\BookSplit{})\BookSplit{}}中只能用+-*/和()运算符，并且只能做整数运算。


\subsection{转义字符\textbackslash{}}
\BookAnchor{db.ch31s03.id2873001}%


和C语言类似，\textbackslash{}在Shell中被用作转义字符，用于去除紧跟其后的单个字符的特殊意义（回车除外），换句话说，紧跟其后的字符取字面值。例如：


\BookTerminal{code/ch31s03/10.txt}{8}{6}


比如创建一个文件名为“\$ \$”的文件可以这样：


\BookTerminal{code/ch31s03/11.txt}{3}{1}


还有一个字符虽然不具有特殊含义，但是要用它做文件名也很麻烦，就是-号。如果要创建一个文件名以-号开头的文件，这样是不行的：


\BookTerminal{code/ch31s03/12.txt}{5}{3}


即使加上\textbackslash{}转义也还是报错：


\BookTerminal{code/ch31s03/13.txt}{5}{3}


因为各种UNIX命令都把-号开头的命令行参数当作命令的选项，而不会当作文件名。如果非要处理以-号开头的文件名，可以有两种办法：


\BookTerminal{code/ch31s03/14.txt}{3}{1}


或者


\BookTerminal{code/ch31s03/15.txt}{3}{1}


\textbackslash{}还有一种用法，在\textbackslash{}后敲回车表示续行，Shell并不会立刻执行命令，而是把光标移到下一行，给出一个续行提示符\textgreater{}，等待用户继续输入，最后把所有的续行接到一起当作一个命令执行。例如：


\BookTerminal{code/ch31s03/16.txt}{5}{3}

\subsection{单引号}
\BookAnchor{db.ch31s03.id2873083}%


和C语言不一样，Shell脚本中的单引号和双引号一样都是字符串的界定符（双引号下一节介绍），而不是字符的界定符。单引号用于保持引号内所有字符的字面值，即使引号内的\textbackslash{}和回车也不例外，但是字符串中不能出现单引号。如果引号没有配对就输入回车，Shell会给出续行提示符，要求用户把引号配上对。例如：


\BookTerminal{code/ch31s03/17.txt}{8}{6}

\subsection{双引号}
\BookAnchor{db.ch31s03.id2873112}%


双引号用于保持引号内所有字符的字面值（回车也不例外），但以下情况除外：


\begin{itemize}

\item \$加变量名可以取变量的值

\item 反引号仍表示命令替换

\item \textbackslash{}\$表示\$的字面值

\item \textbackslash{}`表示`的字面值

\item \textbackslash{}\BookQuoteChar{}表示\BookQuoteChar{}的字面值

\item \textbackslash{}\textbackslash{}表示\textbackslash{}的字面值

\item 除以上情况之外，在其它字符前面的\textbackslash{}无特殊含义，只表示字面值
\end{itemize}

\BookTerminal{code/ch31s03/18.txt}{9}{12}

\BookNav{\hyperlink{db.ch31s02}{上一页}\enspace\textperiodcentered\enspace \hyperlink{db.ch31}{上一级}\enspace\textperiodcentered\enspace \hyperlink{db.index}{起始页}\enspace\textperiodcentered\enspace \hyperlink{db.ch31s04}{下一页}}


% https://akaedu.github.io/book/ch31s04.html

\section{bash启动脚本}
\BookAnchor{db.ch31s04}\BookAnchor{db.ch31s04.id2873178}%
启动脚本是\BookCode{bash}启动时自动执行的脚本。用户可以把一些环境变量的设置和\BookCode{alias}、\BookCode{umask}设置放在启动脚本中，这样每次启动Shell时这些设置都自动生效。思考一下，\BookCode{bash}在执行启动脚本时是以\BookCode{fork}子Shell方式执行的还是以\BookCode{source}方式执行的？



启动bash的方法不同，执行启动脚本的步骤也不相同，具体可分为以下几种情况。


\subsection{作为交互登录Shell启动，或者使用--login参数启动}
\BookAnchor{db.ch31s04.id2873231}%


交互Shell是指用户在提示符下输命令的Shell而非执行脚本的Shell，登录Shell就是在输入用户名和密码登录后得到的Shell，比如从字符终端登录或者用\BookCode{telnet}/\BookCode{ssh}从远程登录，但是从图形界面的窗口管理器登录之后会显示桌面而不会产生登录Shell（也不会执行启动脚本），在图形界面下打开终端窗口得到的Shell也不是登录Shell。



这样启动\BookCode{bash}会自动执行以下脚本：


\begin{enumerate}

\item 首先执行\BookCode{/\BookSplit{}etc/\BookSplit{}profile}，系统中每个用户登录时都要执行这个脚本，如果系统管理员希望某个设置对所有用户都生效，可以写在这个脚本里

\item 然后依次查找当前用户主目录的\BookCode{\textasciitilde{}\BookSplit{}/\BookSplit{}.\BookSplit{}bash\_\BookSplit{}profile}、\BookCode{\textasciitilde{}\BookSplit{}/\BookSplit{}.\BookSplit{}bash\_\BookSplit{}login}和\BookCode{\textasciitilde{}\BookSplit{}/\BookSplit{}.\BookSplit{}profile}三个文件，找到第一个存在并且可读的文件来执行，如果希望某个设置只对当前用户生效，可以写在这个脚本里，由于这个脚本在\BookCode{/\BookSplit{}etc/\BookSplit{}profile}之后执行，\BookCode{/\BookSplit{}etc/\BookSplit{}profile}设置的一些环境变量的值在这个脚本中可以修改，也就是说，当前用户的设置可以覆盖（Override）\BookAnchor{db.ch31s04.id2873327}系统中全局的设置。\BookCode{\textasciitilde{}\BookSplit{}/\BookSplit{}.\BookSplit{}profile}这个启动脚本是\BookCode{sh}规定的，\BookCode{bash}规定首先查找以\BookCode{\textasciitilde{}\BookSplit{}/\BookSplit{}.\BookSplit{}bash\_\BookSplit{}}开头的启动脚本，如果没有则执行\BookCode{\textasciitilde{}\BookSplit{}/\BookSplit{}.\BookSplit{}profile}，是为了和\BookCode{sh}保持一致。

\item 顺便一提，在退出登录时会执行\BookCode{\textasciitilde{}\BookSplit{}/\BookSplit{}.\BookSplit{}bash\_\BookSplit{}logout}脚本（如果它存在的话）。
\end{enumerate}

\subsection{以交互非登录Shell启动}
\BookAnchor{db.ch31s04.id2873387}%


比如在图形界面下开一个终端窗口，或者在登录Shell提示符下再输入\BookCode{bash}命令，就得到一个交互非登录的Shell，这种Shell在启动时自动执行\BookCode{\textasciitilde{}\BookSplit{}/\BookSplit{}.\BookSplit{}bashrc}脚本。



为了使登录Shell也能自动执行\BookCode{\textasciitilde{}\BookSplit{}/\BookSplit{}.\BookSplit{}bashrc}，通常在\BookCode{\textasciitilde{}\BookSplit{}/\BookSplit{}.\BookSplit{}bash\_\BookSplit{}profile}中调用\BookCode{\textasciitilde{}\BookSplit{}/\BookSplit{}.\BookSplit{}bashrc}：


\BookSource{numbers=none}{code/ch31s04/01.pv}{5}{3}


这几行的意思是如果\BookCode{\textasciitilde{}\BookSplit{}/\BookSplit{}.\BookSplit{}bashrc}文件存在则\BookCode{source}它。多数Linux发行版在创建帐户时会自动创建\BookCode{\textasciitilde{}\BookSplit{}/\BookSplit{}.\BookSplit{}bash\_\BookSplit{}profile}和\BookCode{\textasciitilde{}\BookSplit{}/\BookSplit{}.\BookSplit{}bashrc}脚本，\BookCode{\textasciitilde{}\BookSplit{}/\BookSplit{}.\BookSplit{}bash\_\BookSplit{}profile}中通常都有上面这几行。所以，如果要在启动脚本中做某些设置，使它在图形终端窗口和字符终端的Shell中都起作用，最好就是在\BookCode{\textasciitilde{}\BookSplit{}/\BookSplit{}.\BookSplit{}bashrc}中设置。



下面做一个实验，在\BookCode{\textasciitilde{}\BookSplit{}/\BookSplit{}.\BookSplit{}bashrc}文件末尾添加一行（如果这个文件不存在就创建它）：


\BookSource{numbers=none}{code/ch31s04/02.pv}{3}{1}


然后关掉终端窗口重新打开，或者从字符终端\BookCode{logout}之后重新登录，现在主目录下的程序应该可以直接输程序名运行而不必输入路径了，例如：


\BookTerminal{code/ch31s04/03.txt}{3}{1}


就可以了，而不必


\BookTerminal{code/ch31s04/04.txt}{3}{1}


为什么登录Shell和非登录Shell的启动脚本要区分开呢？最初的设计是这样考虑的，如果从字符终端或者远程登录，那么登录Shell是该用户的所有其它进程的父进程，也是其它子Shell的父进程，所以环境变量在登录Shell的启动脚本里设置一次就可以自动带到其它非登录Shell里，而Shell的本地变量、函数、\BookCode{alias}等设置没有办法带到子Shell里，需要每次启动非登录Shell时设置一遍，所以就需要有非登录Shell的启动脚本，所以一般来说在\BookCode{\textasciitilde{}\BookSplit{}/\BookSplit{}.\BookSplit{}bash\_\BookSplit{}profile}里设置环境变量，在\BookCode{\textasciitilde{}\BookSplit{}/\BookSplit{}.\BookSplit{}bashrc}里设置本地变量、函数、\BookCode{alias}等。如果你的Linux带有图形系统则不能这样设置，由于从图形界面的窗口管理器登录并不会产生登录Shell，所以环境变量也应该在\BookCode{\textasciitilde{}\BookSplit{}/\BookSplit{}.\BookSplit{}bashrc}里设置。


\subsection{非交互启动}
\BookAnchor{db.ch31s04.id2873571}%


为执行脚本而\BookCode{fork}出来的子Shell是非交互Shell，启动时执行的脚本文件由环境变量\BookCode{BASH\_\BookSplit{}ENV}定义，相当于自动执行以下命令：


\BookSource{numbers=none}{code/ch31s04/05.pv}{3}{1}


如果环境变量\BookCode{BASH\_\BookSplit{}ENV}的值不是空字符串，则把它的值当作启动脚本的文件名，\BookCode{source}这个脚本。


\subsection{以sh命令启动}
\BookAnchor{db.ch31s04.id2873616}%


如果以\BookCode{sh}命令启动\BookCode{bash}，\BookCode{bash}将模拟\BookCode{sh}的行为，以\BookCode{\textasciitilde{}\BookSplit{}/\BookSplit{}.\BookSplit{}bash\_\BookSplit{}}开头的那些启动脚本就不认了。所以，如果作为交互登录Shell启动，或者使用--login参数启动，则依次执行以下脚本：


\begin{enumerate}

\item \BookCode{/\BookSplit{}etc/\BookSplit{}profile}

\item \BookCode{\textasciitilde{}\BookSplit{}/\BookSplit{}.\BookSplit{}profile}
\end{enumerate}


如果作为交互Shell启动，相当于自动执行以下命令：


\BookSource{numbers=none}{code/ch31s04/06.pv}{3}{1}


如果作为非交互Shell启动，则不执行任何启动脚本。通常我们写的Shell脚本都以\BookCode{\#! /\BookSplit{}bin/\BookSplit{}sh}开头，都属于这种方式。

\BookNav{\hyperlink{db.ch31s03}{上一页}\enspace\textperiodcentered\enspace \hyperlink{db.ch31}{上一级}\enspace\textperiodcentered\enspace \hyperlink{db.index}{起始页}\enspace\textperiodcentered\enspace \hyperlink{db.ch31s05}{下一页}}


% https://akaedu.github.io/book/ch31s05.html

\section{Shell脚本语法}
\BookAnchor{db.ch31s05}\BookAnchor{db.ch31s05.id2873718}%
\subsection{条件测试：test [}
\BookAnchor{db.ch31s05.id2873722}%


命令\BookCode{test}或\BookCode{[\BookSplit{}}可以测试一个条件是否成立，如果测试结果为真，则该命令的Exit Status为0，如果测试结果为假，则命令的Exit Status为1（注意与C语言的逻辑表示正好相反）。例如测试两个数的大小关系：


\BookTerminal{code/ch31s05/01.txt}{9}{10}


\emph{虽然看起来很奇怪，但左方括号\BookCode{[\BookSplit{}}确实是一个命令的名字，传给命令的各参数之间应该用空格隔开}，比如，\BookCode{\$VAR}、\BookCode{-\BookSplit{}gt}、\BookCode{3}、\BookCode{]\BookSplit{}}是\BookCode{[\BookSplit{}}命令的四个参数，它们之间必须用空格隔开。命令\BookCode{test}或\BookCode{[\BookSplit{}}的参数形式是相同的，只不过\BookCode{test}命令不需要\BookCode{]\BookSplit{}}参数。以\BookCode{[\BookSplit{}}命令为例，常见的测试命令如下表所示：


\begingroup
\setlength{\tabcolsep}{6pt}
\small
\sloppy
\renewcommand{\arraystretch}{1.16}
\setlength{\parskip}{0pt}
\begin{longtable}{L{\dimexpr 0.327377\linewidth-2\tabcolsep\relax}L{\dimexpr 0.672623\linewidth-2\tabcolsep\relax}}
\BookTableTitle{测试命令}\BookTableLabel{\BookAnchor{db.ch31s05.id2873825}}
\toprule
\BookCode{[\BookSplit{} -\BookSplit{}d DIR ]\BookSplit{}} & 如果\BookCode{DIR}存在并且是一个目录则为真 \\
\addlinespace[2pt]
\BookCode{[\BookSplit{} -\BookSplit{}f FILE ]\BookSplit{}} & 如果\BookCode{FILE}存在且是一个普通文件则为真 \\
\addlinespace[2pt]
\BookCode{[\BookSplit{} -\BookSplit{}z STRING ]\BookSplit{}} & 如果\BookCode{STRING}的长度为零则为真 \\
\addlinespace[2pt]
\BookCode{[\BookSplit{} -\BookSplit{}n STRING ]\BookSplit{}} & 如果\BookCode{STRING}的长度非零则为真 \\
\addlinespace[2pt]
\BookCode{[\BookSplit{} STRING1 =\BookSplit{} STRING2 ]\BookSplit{}} & 如果两个字符串相同则为真 \\
\addlinespace[2pt]
\BookCode{[\BookSplit{} STRING1 !=\BookSplit{} STRING2 ]\BookSplit{}} & 如果字符串不相同则为真 \\
\addlinespace[2pt]
\BookCode{[\BookSplit{} ARG1 OP ARG2 ]\BookSplit{}} & \BookCode{ARG1}和\BookCode{ARG2}应该是整数或者取值为整数的变量，\BookCode{OP}是\BookCode{-\BookSplit{}eq}（等于）\BookCode{-\BookSplit{}ne}（不等于）\BookCode{-\BookSplit{}lt}（小于）\BookCode{-\BookSplit{}le}（小于等于）\BookCode{-\BookSplit{}gt}（大于）\BookCode{-\BookSplit{}ge}（大于等于）之中的一个 \\
\bottomrule
\end{longtable}
\endgroup


和C语言类似，测试条件之间还可以做与、或、非逻辑运算：


\begingroup
\setlength{\tabcolsep}{6pt}
\small
\sloppy
\renewcommand{\arraystretch}{1.16}
\setlength{\parskip}{0pt}
\begin{longtable}{L{\dimexpr 0.282192\linewidth-2\tabcolsep\relax}L{\dimexpr 0.717808\linewidth-2\tabcolsep\relax}}
\BookTableTitle{带与、或、非的测试命令}\BookTableLabel{\BookAnchor{db.ch31s05.id2873997}}
\toprule
\BookCode{[\BookSplit{} ! EXPR ]\BookSplit{}} & \BookCode{EXPR}可以是上表中的任意一种测试条件，!表示逻辑反 \\
\addlinespace[2pt]
\BookCode{[\BookSplit{} EXPR1 -\BookSplit{}a EXPR2 ]\BookSplit{}} & \BookCode{EXPR1}和\BookCode{EXPR2}可以是上表中的任意一种测试条件，\BookCode{-\BookSplit{}a}表示逻辑与 \\
\addlinespace[2pt]
\BookCode{[\BookSplit{} EXPR1 -\BookSplit{}o EXPR2 ]\BookSplit{}} & \BookCode{EXPR1}和\BookCode{EXPR2}可以是上表中的任意一种测试条件，\BookCode{-\BookSplit{}o}表示逻辑或 \\
\bottomrule
\end{longtable}
\endgroup


例如：


\BookTerminal{code/ch31s05/02.txt}{6}{4}


注意，如果上例中的\BookCode{\$VAR}变量事先没有定义，则被Shell展开为空字符串，会造成测试条件的语法错误（展开为\BookCode{[\BookSplit{} -\BookSplit{}d Desktop -\BookSplit{}a =\BookSplit{} 'abc' ]\BookSplit{}}），作为一种好的Shell编程习惯，应该总是把变量取值放在双引号之中（展开为\BookCode{[\BookSplit{} -\BookSplit{}d Desktop -\BookSplit{}a \BookQuoteChar{}\BookQuoteChar{} =\BookSplit{} 'abc' ]\BookSplit{}}）：


\BookTerminal{code/ch31s05/03.txt}{8}{6}

\subsection{if/then/elif/else/fi}
\BookAnchor{db.ch31s05.id2874121}%


和C语言类似，在Shell中用\BookCode{if}、\BookCode{then}、\BookCode{elif}、\BookCode{else}、\BookCode{fi}这几条命令实现分支控制。这种流程控制语句本质上也是由若干条Shell命令组成的，例如先前讲过的


\BookSource{numbers=none}{code/ch31s05/04.pv}{5}{3}


其实是三条命令，\BookCode{if [\BookSplit{} -\BookSplit{}f \textasciitilde{}\BookSplit{}/\BookSplit{}.\BookSplit{}bashrc ]\BookSplit{}}是第一条，\BookCode{then .\BookSplit{} \textasciitilde{}\BookSplit{}/\BookSplit{}.\BookSplit{}bashrc}是第二条，\BookCode{fi}是第三条。如果两条命令写在同一行则需要用;号隔开，一行只写一条命令就不需要写;号了，另外，\BookCode{then}后面有换行，但这条命令没写完，Shell会自动续行，把下一行接在\BookCode{then}后面当作一条命令处理。和\BookCode{[\BookSplit{}}命令一样，要注意命令和各参数之间必须用空格隔开。\BookCode{if}命令的参数组成一条子命令，如果该子命令的Exit Status为0（表示真），则执行\BookCode{then}后面的子命令，如果Exit Status非0（表示假），则执行\BookCode{elif}、\BookCode{else}或者\BookCode{fi}后面的子命令。\BookCode{if}后面的子命令通常是测试命令，但也可以是其它命令。Shell脚本没有\{\}括号，所以用\BookCode{fi}表示\BookCode{if}语句块的结束。见下例：


\BookSource{numbers=none}{code/ch31s05/05.pv}{9}{7}


\BookCode{:\BookSplit{}}是一个特殊的命令，称为空命令，该命令不做任何事，但Exit Status总是真。此外，也可以执行\BookCode{/\BookSplit{}bin/\BookSplit{}true}或\BookCode{/\BookSplit{}bin/\BookSplit{}false}得到真或假的Exit Status。再看一个例子：


\BookSource{numbers=left,xleftmargin=19pt}{code/ch31s05/06.pv}{9}{13}


上例中的\BookCode{read}命令的作用是等待用户输入一行字符串，将该字符串存到一个Shell变量中。



此外，Shell还提供了\&\&和||语法，和C语言类似，具有Short-circuit特性，很多Shell脚本喜欢写成这样：


\BookPlain{code/ch31s05/07.txt}{3}{1}


\&\&相当于“if.\BookSplit{}.\BookSplit{}.\BookSplit{}then.\BookSplit{}.\BookSplit{}.\BookSplit{}”，而||相当于“if not.\BookSplit{}.\BookSplit{}.\BookSplit{}then.\BookSplit{}.\BookSplit{}.\BookSplit{}”。\&\&和||用于连接两个命令，而上面讲的\BookCode{-\BookSplit{}a}和\BookCode{-\BookSplit{}o}仅用于在测试表达式中连接两个测试条件，要注意它们的区别，例如，


\BookPlain{code/ch31s05/08.txt}{3}{1}


和以下写法是等价的


\BookPlain{code/ch31s05/09.txt}{3}{1}

\subsection{case/esac}
\BookAnchor{db.ch31s05.id2874366}%


\BookCode{case}命令可类比C语言的\BookCode{switch}/\BookCode{case}语句，\BookCode{esac}表示\BookCode{case}语句块的结束。C语言的\BookCode{case}只能匹配整型或字符型常量表达式，而Shell脚本的\BookCode{case}可以匹配字符串和Wildcard，每个匹配分支可以有若干条命令，末尾必须以;;结束，执行时找到第一个匹配的分支并执行相应的命令，然后直接跳到\BookCode{esac}之后，不需要像C语言一样用\BookCode{break}跳出。


\BookSource{numbers=left,xleftmargin=19pt}{code/ch31s05/10.pv}{9}{14}


使用\BookCode{case}语句的例子可以在系统服务的脚本目录\BookCode{/\BookSplit{}etc/\BookSplit{}init.\BookSplit{}d}中找到。这个目录下的脚本大多具有这种形式（以\BookCode{/\BookSplit{}etc/\BookSplit{}apache2}为例）：


\BookPlain{code/ch31s05/11.txt}{9}{17}


启动\BookCode{apache2}服务的命令是


\BookTerminal{code/ch31s05/12.txt}{3}{1}


\BookCode{\$1}是一个特殊变量，在执行脚本时自动取值为第一个命令行参数，也就是\BookCode{start}，所以进入\BookCode{start)\BookSplit{}}分支执行相关的命令。同理，命令行参数指定为\BookCode{stop}、\BookCode{reload}或\BookCode{restart}可以进入其它分支执行停止服务、重新加载配置文件或重新启动服务的相关命令。


\subsection{for/do/done}
\BookAnchor{db.ch31s05.id2874526}%


Shell脚本的\BookCode{for}循环结构和C语言很不一样，它类似于某些编程语言的\BookCode{foreach}循环。例如：


\BookSource{numbers=none}{code/ch31s05/13.pv}{7}{5}


\BookCode{FRUIT}是一个循环变量，第一次循环\BookCode{\$FRUIT}的取值是\BookCode{apple}，第二次取值是\BookCode{banana}，第三次取值是\BookCode{pear}。再比如，要将当前目录下的\BookCode{chap0}、\BookCode{chap1}、\BookCode{chap2}等文件名改为\BookCode{chap0\textasciitilde{}\BookSplit{}}、\BookCode{chap1\textasciitilde{}\BookSplit{}}、\BookCode{chap2\textasciitilde{}\BookSplit{}}等（按惯例，末尾有\textasciitilde{}字符的文件名表示临时文件），这个命令可以这样写：


\BookTerminal{code/ch31s05/14.txt}{3}{1}


也可以这样写：


\BookTerminal{code/ch31s05/15.txt}{3}{1}

\subsection{while/do/done}
\BookAnchor{db.ch31s05.id2874637}%


\BookCode{while}的用法和C语言类似。比如一个验证密码的脚本：


\BookSource{numbers=none}{code/ch31s05/16.pv}{9}{8}


下面的例子通过算术运算控制循环的次数：


\BookSource{numbers=none}{code/ch31s05/17.pv}{9}{7}


Shell还有until循环，类似C语言的do.\BookSplit{}.\BookSplit{}.\BookSplit{}while循环。本章从略。


\subsubsection*{习题}
\BookAnchor{db.ch31s05.id2874672}%


1、把上面验证密码的程序修改一下，如果用户输错五次密码就报错退出。


\subsection{位置参数和特殊变量}
\BookAnchor{db.ch31s05.id2874685}%


有很多特殊变量是被Shell自动赋值的，我们已经遇到了\BookCode{\$?}和\BookCode{\$1}，现在总结一下：


\begingroup
\setlength{\tabcolsep}{6pt}
\small
\sloppy
\renewcommand{\arraystretch}{1.16}
\setlength{\parskip}{0pt}
\begin{longtable}{L{\dimexpr 0.149587\linewidth-2\tabcolsep\relax}L{\dimexpr 0.850413\linewidth-2\tabcolsep\relax}}
\BookTableTitle{常用的位置参数和特殊变量}\BookTableLabel{\BookAnchor{db.ch31s05.id2874706}}
\toprule
\BookCode{\$0} & 相当于C语言\BookCode{main}函数的\BookCode{argv[\BookSplit{}0]\BookSplit{}} \\
\addlinespace[2pt]
\BookCode{\$1}、\BookCode{\$2}... & 这些称为位置参数（Positional Parameter）\BookAnchor{db.ch31s05.id2874755}，相当于C语言\BookCode{main}函数的\BookCode{argv[\BookSplit{}1]\BookSplit{}}、\BookCode{argv[\BookSplit{}2]\BookSplit{}}... \\
\addlinespace[2pt]
\BookCode{\$\#} & 相当于C语言\BookCode{main}函数的\BookCode{argc -\BookSplit{} 1}，注意这里的\BookCode{\#}后面不表示注释 \\
\addlinespace[2pt]
\BookCode{\$@\BookSplit{}} & 表示参数列表\BookCode{\BookQuoteChar{}\$1\BookQuoteChar{} \BookQuoteChar{}\$2\BookQuoteChar{} .\BookSplit{}.\BookSplit{}.\BookSplit{}}，例如可以用在\BookCode{for}循环中的\BookCode{in}后面。 \\
\addlinespace[2pt]
\BookCode{\$?} & 上一条命令的Exit Status \\
\addlinespace[2pt]
\BookCode{\$\$} & 当前Shell的进程号 \\
\bottomrule
\end{longtable}
\endgroup


位置参数可以用\BookCode{shift}命令左移。比如\BookCode{shift 3}表示原来的\BookCode{\$4}现在变成\BookCode{\$1}，原来的\BookCode{\$5}现在变成\BookCode{\$2}等等，原来的\BookCode{\$1}、\BookCode{\$2}、\BookCode{\$3}丢弃，\BookCode{\$0}不移动。不带参数的\BookCode{shift}命令相当于\BookCode{shift 1}。例如：


\BookSource{numbers=none}{code/ch31s05/18.pv}{9}{10}

\subsection{函数}
\BookAnchor{db.ch31s05.id2874943}%


和C语言类似，Shell中也有函数的概念，但是函数定义中没有返回值也没有参数列表。例如：


\BookSource{numbers=none}{code/ch31s05/19.pv}{8}{6}


注意函数体的左花括号\{和后面的命令之间必须有空格或换行，如果将最后一条命令和右花括号\BookCode{\}}写在同一行，命令末尾必须有;号。



在定义\BookCode{foo(\BookSplit{})\BookSplit{}}函数时并不执行函数体中的命令，就像定义变量一样，只是给\BookCode{foo}这个名字一个定义，到后面调用\BookCode{foo}函数的时候（注意Shell中的函数调用不写括号）才执行函数体中的命令。Shell脚本中的函数必须先定义后调用，一般把函数定义都写在脚本的前面，把函数调用和其它命令写在脚本的最后（类似C语言中的\BookCode{main}函数，这才是整个脚本实际开始执行命令的地方）。



Shell函数没有参数列表并不表示不能传参数，事实上，函数就像是迷你脚本，调用函数时可以传任意个参数，在函数内同样是用\BookCode{\$0}、\BookCode{\$1}、\BookCode{\$2}等变量来提取参数，函数中的位置参数相当于函数的局部变量，改变这些变量并不会影响函数外面的\BookCode{\$0}、\BookCode{\$1}、\BookCode{\$2}等变量。函数中可以用\BookCode{return}命令返回，如果\BookCode{return}后面跟一个数字则表示函数的Exit Status。



下面这个脚本可以一次创建多个目录，各目录名通过命令行参数传入，脚本逐个测试各目录是否存在，如果目录不存在，首先打印信息然后试着创建该目录。


\BookSource{numbers=left,xleftmargin=19pt}{code/ch31s05/20.pv}{9}{24}


注意\BookCode{is\_\BookSplit{}directory(\BookSplit{})\BookSplit{}}返回0表示真返回1表示假。

\BookNav{\hyperlink{db.ch31s04}{上一页}\enspace\textperiodcentered\enspace \hyperlink{db.ch31}{上一级}\enspace\textperiodcentered\enspace \hyperlink{db.index}{起始页}\enspace\textperiodcentered\enspace \hyperlink{db.ch31s06}{下一页}}


% https://akaedu.github.io/book/ch31s06.html

\section{Shell脚本的调试方法}
\BookAnchor{db.ch31s06}\BookAnchor{db.ch31s06.id2875095}%
Shell提供了一些用于调试脚本的选项，如下所示：


\begin{description}
\item[{-n}]


读一遍脚本中的命令但不执行，用于检查脚本中的语法错误

\item[{-v}]


一边执行脚本，一边将执行过的脚本命令打印到标准错误输出

\item[{-x}]


提供跟踪执行信息，将执行的每一条命令和结果依次打印出来


\end{description}


使用这些选项有三种方法，一是在命令行提供参数


\BookTerminal{code/ch31s06/01.txt}{3}{1}


二是在脚本开头提供参数


\BookSource{numbers=none}{code/ch31s06/02.pv}{3}{1}


第三种方法是在脚本中用set命令启用或禁用参数


\BookSource{numbers=none}{code/ch31s06/03.pv}{9}{7}


\BookCode{set -\BookSplit{}x}和\BookCode{set +\BookSplit{}x}分别表示启用和禁用\BookCode{-\BookSplit{}x}参数，这样可以只对脚本中的某一段进行跟踪调试。

\BookNav{\hyperlink{db.ch31s05}{上一页}\enspace\textperiodcentered\enspace \hyperlink{db.ch31}{上一级}\enspace\textperiodcentered\enspace \hyperlink{db.index}{起始页}\enspace\textperiodcentered\enspace \hyperlink{db.ch32}{下一页}}
